\SourceLeaf{323} den verlangten mit 0. Ich bezeichne ferner die Neigung der Geraden 12 gegen eine feste $L$ schlechthin mit 12; diese Neigung wird in einem \textls[120]{bestimmten} Sinn als wachsend betrachtet von 0 bis $360^0$ (die Vorstellung zu fixiren denken Sie, $L$ sei der Meridian, und 12 das Azimuth von 2 aus 1 gesehen). Auf ähnliche Weise und auf dieselbe $L$ bezogen, sollen 21, 13, 31, 23, 32 verstanden werden, wo also sich von selbst versteht, dass $21-12=\pm180^0$ (oder allgemeiner einem ungeraden Vielfachen von $180^0$ gleich). Die Zeichen 01, 02, 03 sollen zwar ähnliche Bedeutung haben, jedoch mit dem Unterschiede, dass sie nicht auf $L$, sondern auf \textls[120]{eine andere ganz willkürliche} feste Linie bezogen werden.

Man könnte kurz sagen, 12, 13, 21, 23, 31, 32 sind die wahren Azimuthe der betreffenden Richtungen, hingegen 01, 02, 03 die noch nicht orientirten oder einer noch unbekannten Orientirungscorrection bedürftigen Azimuthe. Es ist unnöthig, zu bemerken, dass hier Azimuth und Meridian bloss zur Fixirung der Vorstellung dienen.

Das Criterium, ob die vorgeschriebenen 6 Grössen (wenn Sie 12 und 21 bloss für Eine zählen) physisch verträglich sind, ist nun dieses:
\[
\left.
\begin{array}{l}
\sin(12-13-02+03)\\
\sin(23-21-03+01)\\
\sin(31-32-01+02)
\end{array}
\right\}\ \text{müssen Einerlei Zeichen haben.}
\]

Das vorhin erwähnte, meines Wissens\HistoricalNote{*)}{*)}{Sollte ich darin mich irren, und Sie mir nachweisen können, dass das Theorem schon von sonst jemand aufgestellt sei, so werden Sie mich sehr verpflichten.} bisher unbekannte, in meinem Besitz seit fast einem halben Jahrhundert befindliche Theorem, die Vierecke betreffend, ist folgendes:
\[
\frac{\sin(12-13-02+03)}{(01).(23)}=\frac{\sin(23-21-03+01)}{(02).(13)}=\frac{\sin(31-32-01+02)}{(03).(12)}\,,
\]

wo die Factoren in den Nennern die durchweg positiv genommenen Abstände zwischen den betreffenden Punkten ausdrücken. Die \textls[120]{Ableitung} dieses schönen Theorems vermittelst des Gebrauchs der complexen Zahlen hat die höchste Einfachheit. Hier will ich sie bloss andeuten.



